\section{Critical assignments and restrictions for pigeonhole resolution}
\label{sec:p07-haken-critical-restrictions}

The exponential resolution result in
Theorem~\ref{thm:p07-haken-resolution-lower-bound} has a fixed
propositional encoding and a complete counting proof in
Annex~\ref{annex:p7.complexity.1}. The following complete elementary
constructions distinguish critical assignments, restrictions and
derivation measures. Their conditional charging inequality isolates
the extra combinatorial estimate that an exponential bound requires
\cite{Haken1985,CookReckhow1979}.

\subsection{The exact clause family and counted assignments}

\begin{hbconvention}{def:p07-haken-host-encoding}{The nonfunctional pigeonhole CNF}
Use the host family \(\mathsf{PHP}^{n}_{n+1}\), \(n\geq1\),
of Definition~\ref{def:p07-pigeonhole-cnf}: variables
\(x_{ij}\) have \(1\leq i\leq n+1\), \(1\leq j\leq n\).
Its row clauses are \(A_i=\{x_{i1},\ldots,x_{in}\}\), and
its collision clauses are
\(B_{ikj}=\{\neg x_{ij},\neg x_{kj}\}\) for i<k.
Clauses are finite sets of literals interpreted disjunctively;
the family is their conjunction. It has
\(n(n+1)\) variables and
\((n+1)+n\binom{n+1}{2}\) input clauses.
No within-pigeon exclusivity clauses are assumed.

The original Cook--Reckhow indexing in
Definition~\ref{def:p07-cr-source-php} has n pigeons and n-1
holes; this host instance is its source \(\mathcal P_{n+1}\)
after atom renaming. A lower bound for an encoding with extra
functionality clauses cannot be transferred here simply by
calling both formulas ``the pigeonhole principle.''
\end{hbconvention}

\begin{hbdefinition}{def:p07-haken-critical}{Critical full truth assignments}
Choose a unique unmatched pigeon \(u\in\{1,\ldots,n+1\}\)
and a bijection \(h:\{1,\ldots,n+1\}\setminus\{u\}
\to\{1,\ldots,n\}\). The associated full truth assignment
\(v_{u,h}\) makes \(x_{ij}\) true exactly when i is not u
and h(i)=j, and makes every other variable false.
Call this assignment \emph{critical}, and write \(\mathcal C_n\)
for this finite family. These assignments are not satisfying
assignments of the whole CNF: each has one specified false row.
Their injection structure is a counting convention, not an added
functionality axiom of the formula.
\end{hbdefinition}

\begin{hbtheorem}{thm:p07-haken-critical-count-routing}{proved=annex:p7.haken.critical-count-routing}{Exact critical count and unique-input routing}
\(\mathcal C_n\) has \((n+1)n!\) assignments. Each satisfies
every collision clause and every row clause other than \(A_u\),
which it falsifies. In any fixed annotated resolution refutation,
start at the empty clause and, at each falsified resolvent, follow
the parent selected by the assignment's value on that inference's
pivot. This gives a terminating route ending at the unique false
input \(A_u\), even when the proof is a shared DAG.
\end{hbtheorem}

\subsection{Matching restrictions and surviving resolution proofs}

\begin{hbdefinition}{def:p07-haken-matching-restriction}{A restriction induced by a partial injection}
Let \(I\subseteq\{1,\ldots,n+1\}\) have r members,
\(0\leq r\leq n\), and let
\(f:I\to\{1,\ldots,n\}\) be injective, with image J.
The partial truth assignment \(\rho_f\) sets every assigned
row i in I to its unique true entry \(x_{i,f(i)}\), and sets
every variable in an unassigned row and occupied column j in J
to false. Variables in unassigned rows and unoccupied columns
remain unset. Equivalently,
\[
\begin{aligned}
 \rho_f(x_{i,f(i)})&=1&& (i\in I),\\
 \rho_f(x_{ij})&=0&& (i\in I,\ j\ne f(i)),\\
 \rho_f(x_{ij})&=0&& (i\notin I,\ j\in J).
\end{aligned}
\]
Restrict a clause by discarding it when any literal is set true,
and otherwise deleting all literals set false. Write
\(C\upharpoonright\rho_f=\topv\) for the discarded satisfied
case; an empty residual literal set is the empty clause, not
the satisfied marker. Restrict the clause family by retaining
precisely its nondiscarded residual clauses.
\end{hbdefinition}

\begin{hbtheorem}{thm:p07-haken-restricted-family}{proved=annex:p7.haken.restricted-family}{Exact residual CNF and compatible critical assignments}
After renaming remaining rows and columns, the restricted family
is \(\mathsf{PHP}^{n-r}_{n-r+1}\). When r=n, this denotes its
explicit degenerate instance: one remaining pigeon, zero holes
and a single empty row clause. When r=0, the family is unchanged.
The number of original critical assignments extending \(\rho_f\)
is exactly
\[
 (n+1-r)(n-r)!.
\]
Their unset-variable restrictions are bijective with the critical
assignments of that residual family. Thus their fraction among
\(\mathcal C_n\) is
\((n+1-r)(n-r)!/((n+1)n!)\), not a product treating different
pigeon/hole literals as independent Bernoulli choices.
\end{hbtheorem}

\begin{hbtheorem}{thm:p07-haken-resolution-restriction}{proved=annex:p7.haken.resolution-restriction}{Restriction preserves refutations without adding weakening}
A resolution refutation of a clause family F restricts to a
resolution refutation of \(F\upharpoonright\rho\), with no
more clause occurrences, for every partial truth assignment
\(\rho\). This holds for DAG proofs, and for tree proofs
when occurrences are retained as occurrences. Surviving derived
clauses may be replaced by subclauses of their direct restrictions.
The proof includes assigned-pivot and missing-pivot cases and
does not invoke a weakening rule absent from
Sheet~\ref{sheet:resolution-cpl}.
\end{hbtheorem}

\subsection{What the counting step proves, and what its hypothesis must supply}

\begin{hbdefinition}{def:p07-haken-first-charge}{First-hit charging on a fixed proof DAG}
Fix an annotated proof DAG and the critical route of every member
of \(\mathcal C_n\). A designated node set S is a \emph{route
barrier} when every route reaches a node in S. Charge an assignment
to its first such node, in the direction from the empty clause
to input clauses. Its capacity at a node is the number of distinct
critical assignments charged there, not the number of walks
passing through the node after tree unfolding.
\end{hbdefinition}

\begin{hbtheorem}{thm:p07-haken-charging}{proved=annex:p7.haken.charging}{The exact conditional counting inequality}
If S is a route barrier and each of its nodes receives at most b
critical assignments, with \(b>0\), then
\[
 \#S\ \geq\ \left\lceil\frac{(n+1)n!}{b}\right\rceil.
\]
In particular, a proved capacity bound
\(b\leq2^{-\beta n}(n+1)n!\) for an absolute
\(\beta>0\) would give at least \(2^{\beta n}\) nodes.
This implication does not prove that such a barrier or such
an exponentially small capacity exists. Those are separate
combinatorial requirements of the lower-bound argument.
\end{hbtheorem}

The input row nodes alone supply a barrier with capacity n! per
row, yielding only n+1 rows. The empty-clause node is itself a
barrier with full capacity \((n+1)n!\), yielding only one node.
These examples show why a valid route construction or a wide
clause mentioned somewhere in the proof is not the missing
exponential bottleneck estimate. In particular, neither the
number of literal occurrences nor the number of pigeon indices
mentioned in a clause can be silently substituted for a source's
more refined bottleneck condition.

\subsection{A worked shared derivation, six routes and size accounting}

\begin{hbtheorem}{thm:p07-haken-small-routes}{proved=annex:p7.haken.small-routes}{Critical routes in the three-pigeon/two-hole example}
For the complete 19-node derivation in
Annex~\ref{annex:p7.cr.small-php}, put
\(x_i=x_{i1}\), \(y_i=x_{i2}\). Its six critical assignments
follow the exact routes in the reciprocal proof below.
Each terminates at its unmatched pigeon's row, with two
assignments per row. The same proof has 19 DAG nodes and
21 clause occurrences after unfolding; it does not acquire
more distinct critical assignments when those inputs are copied.
\end{hbtheorem}

\begin{hbconvention}{def:p07-haken-proof-measures}{Nodes, occurrences and encoded bits}
For these comparisons, a DAG node is one recorded clause
occurrence with its input tag or two parent indices and pivot.
Equal literal sets derived in two separate nodes remain two
occurrences. A tree counts copies created by unfolding reused
parents. Literal width is a different measure, and bit size
also counts variable identifiers and references. With
\(n(n+1)\) binary-indexed variables, a literal identifier costs
\(O(\log(n+2))\) bits. For N nodes and total W recorded
literal occurrences, a direct annotated encoding uses
\(O(W\log(n+2)+N\log(N+2))\) bits, plus input references
when the family is supplied separately.

The original CNF has \(\Theta(n^3)\) literal occurrences
and \(O(n^3\log(n+2))\) encoded bits. An exponential
lower bound in the hole parameter n is therefore not literally
an exponential bound in the input bit length. Resolution
lower bounds concern the specified rule system and encoding;
the independently added extension systems can have shorter
proofs and do not inherit the same lower bound.
\end{hbconvention}

\paragraph{Source and result boundary.}
The Haken lower bound has a complete editorial counting proof in
Annex~\ref{annex:p7.complexity.1}; the elementary helpers here prove
their separate displayed claims. The Cook--Reckhow proof-system framework, the later
Beame--Pitassi presentation and Kraj\'i\v{c}ek's treatment are
separate sources. No claim about all Frege systems or
\(\mathsf P\ne\mathsf{NP}\) follows from a lower bound for
this one resolution system. The edition's dated Frege status
remains Theorem~\ref{thm:p07-frege-lower-bound-boundary}.
